\section{Proofs of the enclosure lemmas}
\label{sec:g-enclosure-proofs}

The harmonic-flat inequality is now established. To pass to general
ends, we will first construct a smooth strict inner barrier and then
take its minimizing enclosure. We prove here the geometric statements
of Section~\ref{sec:g-enclosures}, which guarantee that this enclosure
has the full area, regularity, and completeness required by the
inequality. The arguments themselves use only the fixed-metric
hypotheses stated there.

Return to the smooth exterior $(X,g)$ with inner boundary $B$ and
the compact filling $O\subset M$ of Section~\ref{sec:g-enclosures}.
The class $\mathcal A$ consists of bounded finite-perimeter sets
containing $O$; $P_g$ counts their full ambient perimeter, and
$a_g(B)$ is the infimum of the areas of full smooth enclosing cuts.
The normal at $B$ points into $X$.

\subsection{The full enclosure minimum}
\label{subsec:g-proof-enclosure}

\begin{proof}[Proof of Lemma~\ref{lem:g-enclosure}]
We first establish existence and the equality of infima without using
frontier regularity.  We then use the local obstacle lemma proved
below to identify the exterior.

\emph{Confinement and existence.}
On the coordinate end put $Y=\nabla r/|\nabla r|_g$.  Its divergence is
the outward mean curvature of a coordinate sphere, so
\eqref{eq:g-end} gives
\begin{equation}
 \operatorname{div}_gY=\frac{n-1}{r}+O(r^{-1-q})>0
 \qquad(r>R_0)
 \label{eq:g-outer-divergence}
\end{equation}
for a sufficiently large $R_0$.  For $E\in\mathcal A$ and almost every
$R>R_0$, the BV trace on $S_R=\{r=R\}$ is defined and the perimeter
measure gives $S_R$ zero mass.  Gauss--Green on the bounded region
$E\cap\{r>R\}$ yields
\[
 \int_{E\cap\{r>R\}}\operatorname{div}_gY\,dV_g
 =\int_{\partial^*E\cap\{r>R\}}g(Y,\nu_E)\,dA_g
   -\int_{S_R}\mathbf1_E\,dA_g.
\]
There is no flux at infinity.  The first boundary integral is no
larger than the removed perimeter, while the second is the area added
by clipping $E$ at $S_R$.  Therefore, writing $E_R=E\setminus\{r>R\}$,
\begin{equation}
 P_g(E_R)-P_g(E)
 \le-\int_{E\cap\{r>R\}}\operatorname{div}_gY\,dV_g\le0.
 \label{eq:g-clipping}
\end{equation}
The inequality is strict if the removed region has positive volume.

Fix $R_1>R_0$.  For each competitor choose a good slicing radius in
$(R_0,R_1)$; the radius need not work for any other competitor.
Clipping confines a minimizing sequence to the same compact
truncation at $R_1$.  The competitor $O$ gives a finite perimeter
bound.  BV compactness and lower semicontinuity on a slightly larger
compact set now give a minimizer.  Containment of $O$ and absence from
the specified exterior tail pass to the $L^1$ limit.  Since clipping
did not change the infimum, this is a global minimizer in
$\mathcal A$, not a minimizer constrained by an artificial outer wall.

For any global minimizer, the strict part of \eqref{eq:g-clipping}
excludes positive volume beyond every good slicing radius above
$R_0$.  Choosing one such radius below a fixed
$R_2\in(R_0,R_1)$ confines every minimizer to the truncation at $R_2$.
The same BV argument shows that any sequence of minimizers has an
$L^1$-convergent subsequence.  Its limit remains admissible, and lower
semicontinuity together with minimality makes its perimeter equal to
the common minimum.  This proves the compactness assertion.  The
minimum is positive: a zero-perimeter indicator would be constant
almost everywhere on connected $M$, whereas every admissible set
contains the nonempty open set $O^\circ$ and misses an end tail.

\emph{Recovery by full smooth cuts.}
A smooth exterior $D$ determines the filled competitor
\[
 E_D=O\cup\bigl(X\setminus\operatorname{int}D\bigr),
 \qquad
 P_g(E_D)=\operatorname{Area}_g(\partial_{\mathrm{man}}D).
\]
The identity includes boundary on $B$.  Consequently the perimeter
infimum is at most $a_g(B)$.

For the reverse inequality, start with any $E\in\mathcal A$.
Choose a smooth vector field supported in a collar of $B$ and equal
there to the normal pointing from $O$ into $X$.  Its flow $\Phi_t$
satisfies $\Phi_t(O)\supset O$ for small positive $t$, with a positive
exterior collar between their boundaries.  Hence
$E_t=\Phi_t(E)$ contains a neighborhood of $O$ up to null sets.
The tangential Jacobian of $\Phi_t$ tends uniformly to one on the
compact perimeter support, so $P_g(E_t)\to P_g(E)$.

Fix $t>0$.  Strict BV approximation, with the constant value one
retained on a smaller neighborhood of $O$ and zero retained outside
a fixed compact set, supplies smooth functions $u_j\in[0,1]$ such that
\[
 u_j\longrightarrow\mathbf1_{E_t}\text{ in }L^1,
 \qquad
 \int_M|du_j|_g\,dV_g\longrightarrow P_g(E_t).
\]
For completeness, smooth in finitely many coordinate charts only
between these two constant regions and use a partition of unity;
the usual strict approximation proof is local on precisely that
region.  Choose $\delta_j\downarrow0$ slowly enough that
$\|u_j-\mathbf1_{E_t}\|_{L^1}/\delta_j\to0$.  Coarea and Sard's
theorem give a regular value $s_j\in(\delta_j,1-\delta_j)$ with
\[
 P_g(F_j)\le
 \frac{\int|du_j|_g\,dV_g}{1-2\delta_j}+o(1),
 \qquad F_j=\{u_j>s_j\}.
\]
Also
$\|\mathbf1_{F_j}-\mathbf1_{E_t}\|_{L^1}
 \le\|u_j-\mathbf1_{E_t}\|_{L^1}/\delta_j$.
Lower semicontinuity thus gives $P_g(F_j)\to P_g(E_t)$.
Each $F_j$ contains a neighborhood of $O$ and has compact smooth
boundary.

Retain the component of $M\setminus\overline{F_j}$ containing the
distant end, and denote its closure by $D_j$.  This amounts to
filling the bounded complementary components.  The boundary of a
compact smooth domain has finitely many components; the operation
selects entire boundary components and introduces none.  Thus $D_j$
is a connected smooth exterior, closed in $X$, and its full intrinsic
boundary lies in $\operatorname{int}X$, with
\[
 a_g(B)\le\operatorname{Area}_g(\partial_{\mathrm{man}}D_j)
 \le P_g(F_j).
\]
Let first $j\to\infty$, then $t\downarrow0$.  This proves
$a_g(B)\le P_g(E)$ for every competitor and hence
\eqref{eq:g-minimum}.  Applied to a minimizer, the same argument and
the lower bound $a_g(B)$ show convergence of the resulting cut areas
to its full perimeter.  It also proves independence of the inner
metric extension.

\emph{The connected exterior.}
Lemma~\ref{lem:g-obstacle}, proved in the next subsection, identifies a closed representative
with frontier $T$, whose singular part has zero
$\mathcal H^{n-1}_g$ measure and whose regular charts have exactly
one filled side and one exterior side.  The complement has one
component containing the distant end.  Any other component $U$ is
bounded, with $\partial U\subset T$.  Since
$\mathcal H^{n-1}_g(T)=P_g(E)<\infty$, the finite-boundary-measure
criterion for perimeter gives $U$ finite perimeter.  On its regular
boundary charts its outer normal is opposite to $\nu_E$.  The singular
part has zero perimeter measure, so the perimeter formula for unions
gives
\[
 P_g(E\cup U)=P_g(E)-P_g(U).
\]
The open set $U$ has positive volume and is bounded; its perimeter is
positive, for otherwise its indicator would be constant on connected
$M$.  Filling it contradicts minimality.  Thus the exterior is
connected.  This argument does not require the singular frontier to
be a smooth manifold.

If $n\le7$, Lemma~\ref{lem:g-obstacle} makes $T$ a compact embedded
$C^{1,1}$ hypersurface.  A smooth field transverse to its continuous
outward normal has a flow collar of $T$.  A small positive flow
displaces $T$ into $\operatorname{int}X$ and enlarges its filled side.
In a smaller collar, smooth a $C^1$ defining function in $C^1$ while
retaining positive derivative along the transverse field.  Its smooth
zero set is a small graph over the displaced frontier.  The graph
displacement extends to a collar isotopy, fixed on $O$ and on the
distant end.  The exterior remains connected, so these are full cuts.
Letting the smoothing error and then the flow time tend to zero gives
the asserted $C^1$ and area convergence.

Finally every bounded finite-perimeter superset of a minimizer still
contains $O$, so global minimality gives the claimed outer
minimization directly.
\end{proof}

\subsection{Local regularity at the obstacle}
\label{subsec:g-proof-obstacle}

We establish the local input used above. Its first step converts the
obstacle constraint into a bounded volume error; its second step shows
that contact tangent cones are planes.

\begin{proof}[Proof of Lemma~\ref{lem:g-obstacle}]
The obstacle first converts constrained minimization to an ordinary
almost-minimization estimate.  Extend the outward unit normal of
$O$ to a smooth field $Z$ with $|Z|_g\le1$ and bounded divergence.
It suffices to work in a collar and to cut off farther away.  For a
local replacement $F$, the set $F\cup O$ is admissible.  Perimeter
submodularity and Gauss--Green give, after canceling unchanged
perimeter outside the replacement region $V$,
\begin{align}
 P_g(E;V)&\le P_g(F;V)+P_g(O;V)-P_g(F\cap O;V)\nonumber\\
 &\le P_g(F;V)+\Lambda\operatorname{Vol}_g(E\mathbin\triangle F),
 \label{eq:g-almost-minimality}
\end{align}
where $\Lambda$ bounds $|\operatorname{div}_gZ|$ on the region.
To check the second inequality without assuming that $O$ has finite
total perimeter, take $Z=\nu_O$ on $\partial O\cap V$ and cut it off
in a larger neighborhood.  Since $F\cap O$ and $O$ agree outside a
compact subset of $V$, Gauss--Green for their indicator difference gives
\[
 P_g(O;V)-P_g(F\cap O;V)
 \le\int_{O\setminus F}\operatorname{div}_gZ\,dV_g
 \le\Lambda\operatorname{Vol}_g(E\mathbin\triangle F).
\]
All the perimeters in this calculation are relative to $V$.

Insertion and deletion of small balls in
\eqref{eq:g-almost-minimality}, followed by the relative
isoperimetric inequality, give the usual two-sided volume-density
bounds and upper and lower perimeter-density bounds.  They identify
the frontier of the closed representative with the support of its
perimeter measure.  The codimension-one almost-minimizer regularity
theorem gives a $C^{1,\alpha}$ regular part and a singular set of
dimension at most $n-8$; the Riemannian version uses the smooth area
integrand in coordinate charts.  The error in a radius-$r$ ball is
$O(r^n)$, hence vanishes after rescaling perimeter by $r^{1-n}$.
Blowups are therefore genuine Euclidean perimeter-minimizing cones.
These are the density and dimension-reduction conclusions of
\cite[Chapters~21, 26, and 28]{Maggi2012}; for the smooth minimizing
case see \cite[Chapter~7, Theorem~5.8]{SimonGMT}.

To see why no contact point is singular, take a tangent filled set
at a point of $\partial O$.  It contains the tangent half-space of
$O$, and its boundary cone is supported in the complementary closed
half-space.  That boundary is a stationary multiplicity-one
minimizing cone.  Write $k=n-1$.  On its spherical link the
nonnegative coordinate height $z$ satisfies, weakly,
$\Delta z=-(k-1)z$.  This identity follows by restricting linear
coordinate functions to a stationary cone.  The link is compact and
stationarity includes its singular set, so testing the identity by
the constant function gives
$0=-(k-1)\int z$.  Since $k-1>0$, the height vanishes.  Thus the
boundary cone is the hyperplane bounding the half-space.  Constancy
and the fact that it is the boundary of a set give multiplicity one;
the two density bounds exclude an empty boundary.  Density-one
regularity for almost-minimizers now gives a single contact graph.

Apply \eqref{eq:g-almost-minimality} to flows of both signs in such a
graph cylinder.  The resulting first-variation functional has a
bounded weak mean curvature.  Its graph function $v$ solves a
uniformly elliptic prescribed-mean-curvature equation with bounded
right-hand side.  The coefficients are smooth functions of position,
$v$, and $Dv$, on the bounded gradient range supplied by
$C^{1,\alpha}$ regularity.  Difference quotients first yield local
$W^{2,2}$ regularity.  In nondivergence form the principal coefficients
are $C^{0,\alpha}$ and the right-hand side is bounded; the local
linear $W^{2,p}$ estimates give $v\in W^{2,p}$ for every finite $p$.
Sobolev embedding then gives $C^{1,\alpha}$ for each $\alpha<1$.
See \cite[Chapter~9]{GilbargTrudinger} for these elliptic estimates.
Off the obstacle, all local replacements are admissible, so the weak
mean curvature is zero; interior elliptic regularity gives
smoothness and minimality on every regular patch.

For the stronger low-dimensional contact conclusion we apply the
smooth-obstacle theorem directly: a pure-perimeter minimizer with a
$C^2$ obstacle in ambient dimension below eight has $C^{1,1}$
frontier, smooth away from contact
\cite[Regularity Theorem~1.3(iii), pp.~368--369]{HuiskenIlmanen2001}.
The smooth metric and obstacle satisfy precisely those local
hypotheses.  Its graph functions are consequently in
$W^{2,\infty}$.  Finally the dimension estimate makes the singular
set $\mathcal H^{n-1}$-null, while the regular frontier agrees with
the reduced boundary up to such a null set.  This proves the area
identity.
\end{proof}

\subsection{Detachment and completeness}
\label{subsec:g-proof-detachment}

It remains to use the strict curvature sign. Adding a thin collar
reduces perimeter wherever that collar is not already filled. This
comparison proves detachment before requiring smoothness of the
minimizing frontier.

\begin{proof}[Proof of Lemma~\ref{lem:g-detachment}]
Compactness and strict negativity give a signed-distance collar
$0\le t\le2\varepsilon$ outside $B$ on which $Y=\partial_t$ obeys
\[
 |Y|_g=1,\qquad \operatorname{div}_gY\le-\beta<0.
\]
For a good slicing level $s\in(0,2\varepsilon)$ set
$O_s=O\cup\{0<t<s\}$ and $W_s=O_s\setminus E$.
Gauss--Green, using a smooth extension of $Y$ across $B$, gives
\[
 \int_{W_s}\operatorname{div}_gY\,dV_g
 =\int_{\{t=s\}}\mathbf1_{E^c}\,dA_g
  -\int_{\partial^*E\cap\{0\le t<s\}}g(Y,\nu_E)\,dA_g.
\]
The second integral includes any original frontier at $t=0$.
The perimeter added by adjoining $O_s$ is the first boundary term;
the removed perimeter is at least the second.  Therefore
\begin{equation}
 P_g(E\cup O_s)-P_g(E)
 \le\int_{W_s}\operatorname{div}_gY\,dV_g
 \le-\beta\operatorname{Vol}_g(W_s).
 \label{eq:g-inner-collar}
\end{equation}
The enlarged set is admissible.  Minimality forces $W_s$ to have
zero volume.  For each minimizer choose one good level
$s\in(\varepsilon,2\varepsilon)$.  The common smaller collar
$0\le t<\varepsilon$ is then filled almost everywhere.  The density
bounds for the closed representative exclude any frontier point in
that collar, so it lies in the interior of $E$.

The frontier is now separated from the obstacle.  Every local
perimeter replacement supported sufficiently near it is admissible,
so it is locally perimeter minimizing there.

To express this local minimality in the language of integral currents,
choose a coordinate ball $U$ with compact closure disjoint from the
obstacle, and orient that chart.  Since $E$ has finite perimeter,
$P_g(E;U)<\infty$.  Write $C=\partial[\![E]\!]$ as a current in
$U$.  If $Z$ is an integral $(n-1)$-cycle compactly
supported in $U$, coning in the ball coordinates fills it by an
integral $n$-current $S$ compactly supported in $U$.  A
top-dimensional integral current has the form $S=[\![k]\!]$, where
$k$ is integer-valued and belongs to $BV(U)$; here $k$ has compact
support and $\partial S=Z$.  Thus
$C+Z=\partial[\![\mathbf 1_E+k]\!]$.  The function
$c(t)=\min\{1,\max\{0,t\}\}$ is $1$-Lipschitz and maps integers
to $\{0,1\}$.  Consequently $c(\mathbf 1_E+k)=\mathbf 1_F$ for a
finite-perimeter set $F$ agreeing with $E$ outside a compact subset
of $U$.  The contraction property of $BV$ variation and set
minimality give
\[
 P_g(E;U)\le P_g(F;U)
 \le |D(\mathbf 1_E+k)|_g(U)
 =\mathbf M_g(C+Z).
\]
Since $P_g(E;U)=\mathbf M_g(C)$, this is the local integral-current
minimizing property used by interior boundary regularity.  The
orientation is needed only in this chart; reversing it changes all
current signs and leaves the mass comparison unchanged.

Applying ambient flows
of both signs gives
\[
 \int_{\partial^*E}\operatorname{div}_{T_x\partial^*E}V\,dA_g=0
\]
for each smooth vector field $V$ supported near $T$.  Thus stationarity
holds across the entire frontier, with no extra first-variation term
at the singular set.  Lemma~\ref{lem:g-obstacle} and interior
minimizing-boundary regularity give the stated dimensional conclusions;
see \cite[Chapter~7, Theorem~5.8]{SimonGMT}.  For the integral-current
formulation and its singular-set dimension reduction, see
\cite{Federer1969,Federer1970}.
Detachment removes contact with the obstacle; the dimension restriction
is what removes interior singularities.

Lemma~\ref{lem:g-enclosure} supplies connectedness and boundedness of
the removed region, so the exterior retains exactly the original
end.  It is closed in the complete ambient manifold.  To check
completeness also for intrinsic Cauchy sequences, choose a subsequence
whose successive intrinsic distances are less than $2^{-j}$ and
join its successive terms by curves in $D$ of lengths less than
$2^{1-j}$.  Their concatenation has finite length and therefore an
ambient endpoint.  Closedness puts the endpoint in $D$.  Appending
it to the concatenated curve gives intrinsic convergence, since the
remaining tail lengths tend to zero.  The original Cauchy sequence
then converges as well.  When $n\le7$, the exterior is in addition a
smooth manifold with boundary $T$.

Every bounded finite-perimeter enlargement of $E$ remains a
competitor containing $O$, proving outer minimization.  In the smooth
case any full cut of $D$ determines precisely such an enlargement;
all its components and portions on $T$ contribute to its perimeter.
The full-area identity follows from \eqref{eq:g-minimum} and the
zero $(n-1)$-dimensional measure of the singular set.
\end{proof}
